% !TeX root = ../main.tex
\section{Proposed Methods for Secondary-Storage KV Offloading}
\label{sec:offloading}

The following subsections organize candidate methods; they do not claim that the mechanisms are implemented or validated.
The evaluated subset uses pinned host DRAM: block offload, on-demand restoration, asynchronous predictive prefetching, inclusive clean copies, two-tier eviction, and fetch-aware scheduling are implemented.
The SSD-specific portions below remain design requirements and are not presented as measured contributions.

\subsection{Architecture and Design Goals}
One radix tree tracks both GPU and host residency.
Each node may hold GPU pages, host pages, or a host-only suffix; GPU-resident nodes remain prefix-closed.
The pinned host pool stores 16-token pages in page-major form and serves as the implemented capacity tier.
CacheWise priority orders GPU and host eviction, while the scheduler restores host matches before prefill.
A future SSD tier would sit below a bounded pinned staging pool, with versioned blocks and explicit read/write queues; the current 37.5~GiB pinned pool is not such a staging design.

\subsection{Method 1: Block-Based Offload and On-Demand Restore}
The implemented baseline operates on 16-token pages.
When GPU allocation needs space, unreferenced radix leaves are split at a page boundary and evicted by the selected priority.
A dirty node receives host pages and a device-to-host copy; a node with a clean inclusive host copy drops its GPU pages without another write.
Prefix lookup returns the deepest GPU handle plus any compatible host suffix, which is locked, allocated new GPU pages, copied back, and reattached before use.

Host locks prevent reclamation during restoration, and host pressure evicts host-only leaves before clean copies of GPU-resident nodes.
If host allocation fails, the subtree is discarded and later recomputed.
Identity is process-local to the radix key and fixed engine configuration.
Persistent SSD blocks would additionally need model and tokenizer revisions, KV dtype/layout, layer geometry, token-prefix digest, positions, length, checksum, and format version; incomplete or incompatible blocks must be invalidated.

\subsection{Method 2: Asynchronous Transfers and Prefetching}
On-demand copies are ordered on the engine CUDA stream; predictive prefetch uses a side stream after the current forward pass, with events protecting later reads and host-slot reuse.
Transfers coalesce contiguous host runs and process at most 256 pages per staging chunk.
Inactive sessions predicted to return within 30~s are considered, with an 8,192-token per-step budget.
Prefetch exchanges GPU pages only when enough victims have predicted return later than $2\tau+1$~s, limiting thrashing.
The prototype has no SSD I/O queue or cancellation; those controls, partial-I/O handling, and bounded staging backpressure remain future work.

\subsection{Method 3: Cost-Aware Offload and Recovery Decisions}
The implemented policy uses CacheWise return-time estimates for eviction and prefetch ordering, not as calibrated reuse probabilities.
For scheduling, request $i$ has cost $c_i=(d_i-k_i)-(1-\alpha)h_i$, where $k_i$ is the GPU prefix, $h_i$ the additional host prefix, and $\alpha=0.05$ the assumed fetch-to-recompute cost ratio.
Thus host-resident tokens are cheaper than missing tokens.
Restoration is used when a host match exists; otherwise the engine recomputes.
An SSD policy must replace fixed $\alpha$ with queue-aware measured cost and admit a write only when expected avoided recomputation covers both write and later read cost.

\subsection{Implementation Scope and Reproducibility}
The artifact is Mini-SGLang commit \texttt{3f9aee7}.\footnote{\url{https://github.com/BoxuanYang/MiniSGLangs}}
The radix cache and host pool implement residency and copies; scheduler cache and prefill modules implement restoration, prefetch, and fetch-aware ordering.
Experiments use \texttt{replay.py} and \texttt{run\_experiments.py}.
Configuration exposes host capacity, policy, predictor, 30-s horizon, 8,192-token prefetch budget, and fetch ratio.
Randomized cache tests and a real-model checker validate restored KV and greedy output.
SSD I/O, durable metadata, a persistent index, direct I/O, page-cache control, bounded staging independent of cache capacity, compression, failure recovery, and cross-process reuse are unimplemented.
